\subsection{Delay-and-Sum Beamformer (DAS)}\label{sec:doa:das}

El método más clásico y sencillo para estimar la dirección de arribo es el 
\textbf{beamformer de retraso y suma} (DAS, por sus siglas en inglés).  
La idea es intuitiva: si se supone un ángulo de llegada, se compensan los retardos 
con los que la onda alcanza a cada sensor, y luego se suman las señales.  

\begin{itemize}
	\item Cuando el ángulo supuesto coincide con el real, las señales quedan alineadas en fase 
	y se refuerzan al sumarse (suma constructiva).  
	\item Si el ángulo no corresponde, las señales llegan desfasadas 
	y se atenúan entre sí (suma destructiva).  
\end{itemize}

De esta forma, basta recorrer un conjunto de ángulos candidatos y calcular la potencia 
de la salida del beamformer para obtener el \textbf{espectro espacial DAS}: 
los máximos de este espectro indican las direcciones estimadas de llegada de la señal 
\cite{YangBook}.

\subsubsection*{Características principales}
\begin{itemize}
	\item Implementación extremadamente simple y de bajo costo computacional.  
	\item No requiere información adicional sobre el entorno ni sobre las señales.  
	\item Resolución angular limitada: los lóbulos principales son anchos, 
	lo que dificulta separar fuentes cercanas.  
	\item Rechazo a interferencias poco eficiente.  
	\item Mejora con mayor número de sensores y mayor apertura del arreglo.  
\end{itemize}

\subsubsection*{Formulación matemática}

Sea $\mathbf{u}(t) \in \mathbb{C}^M$ el vector de observaciones de un arreglo de $M$ sensores, 
y $R_{uu}=\mathbb{E}\{\mathbf{u}(t)\mathbf{u}^H(t)\}$ su matriz de correlación.  
El \emph{steering vector} para un ángulo candidato $\varphi$ es:

\begin{equation}
	\mathbf{a}(\varphi) =
	\begin{bmatrix}
		1 \\
		e^{j 2 \pi d_{\text{norm}} \cos \varphi} \\
		e^{j 2 \pi d_{\text{norm}} \cdot 2 \cos \varphi} \\
		\vdots \\
		e^{j 2 \pi d_{\text{norm}} (M-1)\cos \varphi}
	\end{bmatrix},
\end{equation}

donde $d_{\text{norm}} = d/\lambda$ es el espaciamiento normalizado entre elementos.  

La salida del beamformer apuntando hacia $\varphi$ es:
\begin{equation}
	y(t) = \mathbf{a}^H(\varphi)\, \mathbf{u}(t),
\end{equation}

y su potencia promedio:
\begin{equation}
	P_{DAS}(\varphi) = \mathbf{a}^H(\varphi)\, R_{uu}\, \mathbf{a}(\varphi).
\end{equation}

Al graficar $P_{DAS}(\varphi)$ en función del ángulo, los picos corresponden 
a las direcciones estimadas de arribo.  
En la Fig.~\ref{fig:doa:dasMatlab} se muestra un ejemplo de este espectro 
simulado en MATLAB para un arreglo lineal.

\insertarfigura{fig/DAS.png}{Espectro espacial obtenido con el beamformer Delay-and-Sum (DAS). 
	Se simularon tres fuentes incidentes en $30^\circ$, $50^\circ$ y $70^\circ$, 
	con una relación señal-ruido (SNR) de 5~dB, usando un arreglo de 10 sensores 
	uniformemente espaciados con $d_{norm}=0,5$. 
	Los picos del espectro corresponden a las direcciones estimadas de arribo.}{fig:doa:dasMatlab}{1}
